\section{Lean formalization architecture}\label{sec:lean}

\subsection{Target the existing checked objects}

The implementation should live in the current ProveIt workspace and use the ordinary object-language presentation first.  The essential checked declarations are:

\begin{center}
\begin{tabularx}{\textwidth}{@{}lX@{}}
\toprule
Declaration & Role \\
\midrule
\code{PR k} & fixed primitive-recursive symbols of arity \(k\) \\
\code{PTm n} & terms with \(n\) free numerical variables \\
\code{Fm PTm n} & first-order formulas over those terms \\
\code{PRAAx} & defining equations, successor axioms, and quantifier-free induction \\
\code{Prf PRAAx Gamma phi} & required object-level endpoint \\
\code{PRA.precZero}, \code{PRA.precSucc} & equations for concrete primitive recursions \\
\code{PRA.qfInd} & induction on a syntactically quantifier-free invariant \\
\code{Prf.subst}, \code{Prf.weaken}, \code{Prf.trans} & structural proof composition \\
\code{Prf.allE}, \code{Prf.allI} & specialization and universal closure \\
\code{Prf.eqCongr}, \code{Prf.eqSubst1} & equality transport \\
\code{pra_iff_native} & optional export after an ordinary object proof exists \\
\bottomrule
\end{tabularx}
\end{center}

The checked \code{IPRADerivations} library adds crucial conveniences: \code{inst} for using closed derivations at arbitrary terms, \code{ind} for open quantifier-free induction, concrete arithmetic and order symbols, pairing, finite sums/products, and a proved least bounded search.  Reuse of these modules avoids much of the formalization debt identified by earlier packets.

\subsection{The normalized reflection kernel}

The archive's \code{lean/NormalizedTraceReflection.lean} follows this record shape:

\begin{lstlisting}[language={}] 
structure PRTraceSystem where
  good   : PR 1
  init   : PR 1
  bound  : PR 1
  step   : PR 3       -- counter, state, input
  inv    : PR 3       -- Boolean invariant
  accept : PR 2       -- final state, input

structure NormalizedTraceProofs (T : PRTraceSystem) where
  base : [GoodFm T] |-[PRAAx]
    (InvRunFm T).subst (inst PTm.zero)
  step : InvRunFm T :: [GoodFm T].map Fm.shift |-[PRAAx]
    (InvRunFm T).subst (repl (PTm.succ (var 0)))
  finish : [GoodFm T] |-[PRAAx]
    ((InvRunFm T).subst (inst (boundTerm T)) --> AcceptAtBoundFm T)
\end{lstlisting}

The theorem body is only:

\begin{lstlisting}[language={}] 
have hall : [GoodFm T] |-[PRAAx] Fm.all (InvRunFm T) :=
  PRA.qfInd (InvRunFm T) (invRun_qf T) h.base h.step
have hat := Prf.allE (boundTerm T) hall
have hacc := Prf.impE h.finish hat
exact Prf.allI (Prf.impI hacc)
\end{lstlisting}

Minor \code{simpa} normalization may be required by the exact elaborator, but there is no remaining mathematical idea in the kernel.  This record deliberately accepts normalized obligations.  A separate raw-contract elaborator is responsible for run equations and substitutions.

\subsection{Raw local contracts}

A convenient second record contains readable local formulas:

\begin{lstlisting}[language={}] 
structure RawTraceProofs (T : PRTraceSystem) where
  base  : ([] : List (Fm PTm 1)) |-[PRAAx] BaseFm T
  step  : ([] : List (Fm PTm 3)) |-[PRAAx] StepFm T
  final : ([] : List (Fm PTm 2)) |-[PRAAx] FinalFm T
\end{lstlisting}

The function \code{normalizeRaw} should construct \code{NormalizedTraceProofs}.  Its proof obligations are local and explicit:

\begin{enumerate}
\item instantiate \code{base} at \(u\), eliminate its implication under \code{GoodFm}, and transport the state along \code{run_zero};
\item instantiate \code{step} at \((i,\Run(i,u),u)\), eliminate its implications using the induction and good-input hypotheses, and transport along \code{run_succ};
\item specialize the invariant at the bound, instantiate \code{final} at \((\Run(\mathsf{bound}(u),u),u)\), and form the terminal implication.
\end{enumerate}

The existing checked \code{IPRADerivations.Core.Basic.inst} theorem removes much context/substitution boilerplate.  The remaining term congruence follows the pattern of \code{congrApp}; the logical steps are \code{Prf.impE}, \code{Prf.impI}, \code{Prf.hyp0}, and weakening.  This is the first recommended Lean milestone.

\subsection{Bounded formulas and strategy certificates}

Do not create a universal object-language evaluator for arbitrary formulas.  Instead define a metatheoretic syntax whose atoms and bounds carry concrete \code{PR} expressions, and compile each fixed formula to concrete symbols.

\begin{lstlisting}[language={}] 
inductive BForm (n : Nat)
  | atom   (test : PR n)
  | and    (p q : BForm n)
  | or     (p q : BForm n)
  | exists (bound : PR n) (body : BForm (n+1))
  | forall (bound : PR n) (body : BForm (n+1))
\end{lstlisting}

In an initial implementation, it may be simpler for \code{BForm} to exist only in the host metatheory and for the compiler to emit fixed \code{PR} symbols directly.  Prove checked object lemmas for Boolean \code{allBelow}, \code{anyBelow}, and least search.  Then define certificate encodings and a deterministic checker.  The soundness theorem is generated by structural recursion on the fixed formula, with all object inductions delegated to the fold lemmas.

The negative certificate is essential.  At an existential node, failure checks all children; at a universal node, failure gives one counterexample.  This makes failed proof automation auditable and supports the finite challenge decomposition of \Cref{sec:games}.

\subsection{Relational viability}

Use a first object-level implementation with a rectangular bounded state space:

\begin{lstlisting}[language={}] 
structure PRRelMachine where
  start       : PR 1
  answerBound : PR 3
  halt        : PR 3
  local       : PR 4
  next        : PR 4
  post        : PR 3
  stateBound  : PR 2
\end{lstlisting}

Pair \((r,s,x)\) into the trace state.  The development order is:

\begin{enumerate}
\item define Boolean all/any folds and bounded table lookup/update;
\item define the reverse row recurrence and prove its equation;
\item prove viability correctness by open quantifier-free induction on \(r\);
\item define the least viable answer using checked bounded search;
\item define forward replay and normalize it to a trace;
\item prove transcript soundness and terminal postcondition;
\item add confluence/diameter theorems and export only stable observations.
\end{enumerate}

A sparse external-certificate mode should follow after the rectangular version is checked.  It has a smaller execution footprint but a more elaborate pointer checker.

\subsection{Recommended module graph}

\begin{center}
\begin{tikzpicture}[
 node distance=5mm and 13mm,
 box/.style={draw,rounded corners,align=center,fill=codegray,inner sep=4pt,font=\small},
 arr/.style={-{Latex[length=2mm]},thick}
]
\node[box] (up) {ConciseIPRA\\IPRADerivations Core};
\node[box,below=of up] (codes) {CBIP/Codes\\CBIP/BooleanFold};
\node[box,below left=of codes] (trace) {CBIP/TraceReflection\\CBIP/RawNormalize};
\node[box,below right=of codes] (game) {CBIP/BoundedGame\\CBIP/StrategyCert};
\node[box,below=of trace] (rel) {CBIP/Relational\\CBIP/Stability};
\node[box,below=of game] (err) {CBIP/DemandDAG\\CBIP/ErrorDAG};
\node[box,below=12mm of $(rel)!0.5!(err)$] (obj) {CBIP/RegularNames, FiniteNet, RootCluster, HeightSearch};
\node[box,below=of obj] (apps) {Packet adapters and end-to-end checked theorems};
\draw[arr] (up)--(codes);
\draw[arr] (codes)--(trace);
\draw[arr] (codes)--(game);
\draw[arr] (trace)--(rel);
\draw[arr] (game)--(err);
\draw[arr] (rel)--(obj);
\draw[arr] (err)--(obj);
\draw[arr] (obj)--(apps);
\end{tikzpicture}
\end{center}

\subsection{Dependency-ordered completion plan}

\begin{enumerate}
\item Compile the no-placeholder normalized reflection kernel.
\item Complete \code{normalizeRaw} using \code{inst}, \code{run_zero}, \code{run_succ}, and equality transport.
\item Reuse or promote checked Boolean, pairing, and search modules into a stable public API.
\item Prove \code{allBelow}/\code{anyBelow} elimination and introduction lemmas.
\item Complete one fixed \(\forall\exists\forall\) strategy certificate end to end.
\item Implement rectangular viability and least-viable replay.
\item Implement stable quotient/diameter erasure using the already compiled observable-glue pattern.
\item Add regular-name addition, product, and guarded reciprocal.
\item Add error-DAG and finite-net argmin.
\item Port one substantial packet end to end; good candidates are a separated rotation theorem, a finite-law probability estimate, or a root-cluster symmetric identity.
\item Only after correctness, optimize proof size or target the compressed native calculi.
\end{enumerate}

\subsection{Audits for every exported theorem}

A completed module should pass all of the following:

\begin{itemize}
\item clean \code{lake build};
\item source scan for \code{sorry}, \code{admit}, and newly introduced axioms;
\item \code{#print axioms} on each public theorem;
\item theorem type is visibly \code{Prf PRAAx} before any native export;
\item exact finite regression tests for the associated evaluator/checker;
\item a ledger entry naming every unresolved external producer.
\end{itemize}
